\textbf{Phần 1. Trắc nghiệm lựa chọn}
\setcounterpageref{ex}{0}

\begin{ex} %[2D10V52]
Số $645$ được viết thành tổng các trăm, chục, đơn vị là
\choice
{$600 + 45$}
{$640 + 5$}
{\True $600 + 40 + 5$}
{$6 + 4 + 5$}
\loigiai{
$645 = 600 + 40 + 5$.
\textbf{Đáp án của câu hỏi là C}
}
\end{ex}

\begin{ex} %[2D10V52]
Tổng $800 + 70 + 2$ bằng số nào sau đây?
\choice
{$827$}
{\True $872$}
{$807$}
{$782$}
\loigiai{
$800 + 70 + 2 = 872$.
\textbf{Đáp án của câu hỏi là B}
}
\end{ex}

\begin{ex} %[2D10V52]
Số $509$ được viết thành tổng là
\choice
{$500 + 90$}
{$50 + 9$}
{$500 + 90 + 0$}
{\True $500 + 9$}
\loigiai{
Hàng chục bằng 0 nên $509 = 500 + 9$.
\textbf{Đáp án của câu hỏi là D}
}
\end{ex}

\begin{ex} %[2D10V52]
Cho số có dạng $a37 = 400 + 30 + 7$. Chữ số $a$ là
\choice
{\True $4$}
{$3$}
{$7$}
{$0$}
\loigiai{
Chữ số hàng trăm là 4.
\textbf{Đáp án của câu hỏi là A}
}
\end{ex}

\begin{ex} %[2D10V52]
Số nào dưới đây có chữ số hàng chục bằng $0$ và được viết thành tổng $700 + 8$?
\choice
{$780$}
{$788$}
{\True $708$}
{$807$}
\loigiai{
$700 + 8 = 708$.
\textbf{Đáp án của câu hỏi là C}
}
\end{ex}

\begin{ex} %[2D10V52]
Biết $x = 300 + 60 + 4$ và $y = 300 + 50 + 9$. Khẳng định nào sau đây là đúng?
\choice
{$x < y$}
{\True $x > y$}
{$x = y$}
{Không so sánh được}
\loigiai{
$x = 364, y = 359$. Vì $364 > 359$ nên $x > y$.
\textbf{Đáp án của câu hỏi là B}
}
\end{ex}

\begin{ex} %[2D10V52]
Một số gồm $b$ trăm, $5$ chục và $1$ đơn vị. Biết $b$ là số liền sau số $6$. Số đó là
\choice
{$561$}
{$651$}
{$851$}
{\True $751$}
\loigiai{
Liền sau 6 là 7, nên $b = 7$. Số đó là 751.
\textbf{Đáp án của câu hỏi là D}
}
\end{ex}

\begin{ex} %[2D10V52]
Tổng nào sau đây có giá trị bằng $990$?
\choice
{\True $900 + 90$}
{$900 + 9$}
{$90 + 9$}
{$900 + 90 + 9$}
\loigiai{
$900 + 90 = 990$.
\textbf{Đáp án của câu hỏi là A}
}
\end{ex}

\begin{ex} %[2D10V52]
Cho số $435$. Nếu bớt đi $3$ chục và $5$ đơn vị thì số đó còn lại bao nhiêu?
\choice
{$430$}
{$405$}
{\True $400$}
{$300$}
\loigiai{
$435 - 35 = 400$.
\textbf{Đáp án của câu hỏi là C}
}
\end{ex}

\begin{ex} %[2D10V52]
Số chẵn lớn nhất có ba chữ số được phân tích thành tổng các trăm, chục, đơn vị là
\choice
{$900 + 90 + 9$}
{\True $900 + 90 + 8$}
{$900 + 80 + 8$}
{$800 + 90 + 8$}
\loigiai{
Số chẵn lớn nhất có 3 chữ số là 998. Phân tích: $900 + 90 + 8$.
\textbf{Đáp án của câu hỏi là B}
}
\end{ex}

\begin{ex} %[2D10V52]
Một cửa hàng bán được $200 + 40 + 8$ mét vải. Buổi sáng bán được $100$ mét. Buổi chiều bán được bao nhiêu mét vải?
\choice
{$138\text{ m}$}
{$140\text{ m}$}
{$108\text{ m}$}
{\True $148\text{ m}$}
\loigiai{
Tổng số vải là 248 m. Buổi chiều bán: $248 - 100 = 148$ m.
\textbf{Đáp án của câu hỏi là D}
}
\end{ex}

\begin{ex} %[2D10V52]
Tìm $x$ biết: $x - 70 = 500 + 3$.
\choice
{\True $573$}
{$503$}
{$570$}
{$433$}
\loigiai{
$x - 70 = 503 \Rightarrow x = 503 + 70 = 573$.
\textbf{Đáp án của câu hỏi là A}
}
\end{ex}

\textbf{Phần 2. Trắc nghiệm đúng sai}
\setcounterpageref{ex}{0}

\begin{ex} %[2D10V52]
Xét cấu tạo phân tích số thành tổng trăm, chục, đơn vị:
\choiceTF
{\True Số 320 phân tích thành 300 + 20.}
{Số 406 phân tích thành 400 + 60.}
{\True Biểu thức 900 + 90 + 9 có giá trị bằng số lớn nhất có ba chữ số.}
{\True Số 604 được viết là 600 + 4.}
\loigiai{
\begin{itemize}
    \item Ý a) Đúng: Đúng, hàng đơn vị bằng 0.
    \item Ý b) Sai: Sai, phải là 400 + 6.
    \item Ý c) Đúng: 900 + 90 + 9 = 999.
    \item Ý d) Đúng: Đúng với cấu tạo số có hàng chục bằng 0.
\end{itemize}
\begin{itemize}
    \item a) \textbf{Đúng}.
    \item b) \textbf{Sai}.
    \item c) \textbf{Đúng}.
    \item d) \textbf{Đúng}.
\end{itemize}
}
\end{ex}

\begin{ex} %[2D10V52]
Cho số M = 746:
\choiceTF
{\True M = 700 + 40 + 6.}
{\True Giá trị của chữ số 7 gấp 100 lần giá trị của chữ số 7 ở hàng đơn vị (nếu có).}
{Nếu xóa chữ số 4 thì được số 76.}
{\True Tổng các chữ số của M là 17.}
\loigiai{
\begin{itemize}
    \item Ý a) Đúng: Đúng cấu tạo.
    \item Ý b) Đúng: 700 : 7 = 100. Đúng.
    \item Ý c) Sai: Xóa chữ số 4 thì hàng trăm là 7, đơn vị là 6 thành số 76 (sai, thành 76 thì mất chữ số hàng trăm, thực tế là 76 nếu bỏ hàng chục nhưng theo vị trí là 76). Khẳng định đúng là số còn lại có 2 chữ số.
    \item Ý d) Đúng: 7 + 4 + 6 = 17. Đúng.
\end{itemize}
\begin{itemize}
    \item a) \textbf{Đúng}.
    \item b) \textbf{Đúng}.
    \item c) \textbf{Sai}.
    \item d) \textbf{Đúng}.
\end{itemize}
}
\end{ex}

\textbf{Phần 3. Trả lời ngắn}
\setcounterpageref{ex}{0}

\begin{ex} %[2D10V52]
Tìm số có ba chữ số biết khi xóa chữ số 5 ở hàng trăm thì được số 48.
\shortans{548}
\loigiai{
\textbf{Đáp số:} $548$\\
Hàng trăm là 5, hai hàng sau là 48: số 548.
}
\end{ex}

\begin{ex} %[2D10V52]
Tính giá trị của biểu thức: $A = 700 + 80 + 5 - 200 - 5$.
\shortans{580}
\loigiai{
\textbf{Đáp số:} $580$\\
$A = 785 - 205 = 580$.
}
\end{ex}

\begin{ex} %[2D10V52]
Bác nông dân thu được 6 trăm, 8 chục và 9 bắp ngô. Bác biếu 89 bắp. Còn lại bao nhiêu bắp ngô?
\shortans{600}
\loigiai{
\textbf{Đáp số:} $600$\\
Tổng: 689 bắp. Còn lại: $689 - 89 = 600$ bắp ngô.
}
\end{ex}

\begin{ex} %[2D10V52]
Viết số gồm 9 trăm, 9 chục và 9 đơn vị.
\shortans{999}
\loigiai{
\textbf{Đáp số:} $999$\\
Viết là 999.
}
\end{ex}

\begin{ex} %[2D10V52]
Tìm số tự nhiên biết số đó bằng tổng của 3 trăm và 5 đơn vị.
\shortans{305}
\loigiai{
\textbf{Đáp số:} $305$\\
$300 + 5 = 305$.
}
\end{ex}

\begin{ex} %[2D10V52]
Cho $x = 400 + 20 + 8$. Tính $x - 28$.
\shortans{400}
\loigiai{
\textbf{Đáp số:} $400$\\
$428 - 28 = 400$.
}
\end{ex}
